% Generated draft: requires mathematical and source review.
\documentclass[11pt]{article}
\usepackage[margin=25mm]{geometry}
\usepackage{fontspec}
\setmainfont{Latin Modern Roman}
\usepackage{amsmath,amssymb,mathtools,unicode-math}
\setmathfont{Latin Modern Math}
\usepackage{enumitem,needspace,xurl,hyperref}
\setlength{\parindent}{0pt}
\setlength{\parskip}{5pt}
\setlength{\emergencystretch}{3em}
\newcommand{\N}{\mathbb{N}}
\newcommand{\Z}{\mathbb{Z}}
\newcommand{\Q}{\mathbb{Q}}
\newcommand{\R}{\mathbb{R}}
\newcommand{\C}{\mathbb{C}}
\newcommand{\F}{\mathbb{F}}
\newcommand{\A}{\mathbb{A}}
\newcommand{\PP}{\mathbb{P}}
\newcommand{\E}{\mathbb{E}}
\newcommand{\eps}{\varepsilon}
\newcommand{\dd}{\,\mathrm{d}}
\begin{document}
\section*{UnsolvedMath ID: 11000133}
\textbf{Source problem code:} catalogue entry
\subsection{Definitions and mathematical statement}
Let the two elementary transformations and their coordinate domains be those defined before this problem; their formulas and geometric meanings are not reproduced in the excerpt. The request is to express both transformations usefully by piecewise tropical formulas, and thereby obtain the stated piecewise tropical-polynomial mapping-class-group representation. The formulas must reproduce the original coordinate changes on their admissible domains and respect the move relations needed to descend to the mapping class group; no faithfulness is expressly required.

For orientation about terminology, in a max-plus convention tropical addition and multiplication are \(a\oplus b=\max(a,b)\) and \(a\odot b=a+b\). A tropical polynomial is then a finite maximum of affine-linear forms with nonnegative integral exponent vectors; allowing negative exponents gives Laurent expressions. A min-plus convention is also possible. The source specifies neither convention nor allowed subtraction/Laurent operations or chamber boundaries, so this reference definition does not choose the missing formula class. The chart domains may have constraints and need not be a full Euclidean space.

The source asks what properties the resulting representation has, particularly under iteration of a mapping class. Once its coordinate action \(T_\varphi\) is specified, this means analyzing \(T_\varphi^n\) and its relation to iteration on the geometric objects. No particular growth law, invariant or dynamical classification is asserted in the excerpt.

An alternative suggested coordinate system, attributed to [9], uses intersection numbers with a family of curves \((c_j)_{j\in J}\). For a compactly supported measured foliation \(\mathcal F\), it records
\[I(\mathcal F)=\bigl(i(\mathcal F,c_j)\bigr)_{j\in J},\]
where \(i\) is geometric measured intersection number. The source says these numbers coordinatize the foliations, subject to relations, and tentatively suggests that Theorem 3 could compute them. The curve family, relations and Theorem 3 are absent, so no assertion of effective recovery or unrestricted independent coordinates is inserted. The final request for further applications/consequences is open-ended and supplies no additional fixed theorem.
\subsection{Short English statement}
Give tropical formulas for two specified but omitted coordinate transformations, study the induced mapping-class action under iteration, and explore intersection-number coordinates and further applications. The source's tentative computational suggestion remains tentative.
\subsection{Sources}
\begin{itemize}
\item Local UnsolvedMath registry, permanent ID \(11000133\), exact source code \texttt{catalogue entry}.
\item Farb, \emph{Problems on Mapping Class Groups and Related Topics}, Problem 3, PDF page 117: \url{https://www.math.uchicago.edu/~farb/papers/mcgbook.pdf}. Registry access 2026-07-29; pdf-font-structure extraction.
\item Internal reference [9] and Theorem 3 provide the alternative coordinates and proposed computation, but their content/full citation is not supplied. The summary only mentions unspecified tropical Teichm\"uller literature.
\end{itemize}
\end{document}
